\BourbakiExercisesSection

\begin{enumerate}
\def\labelenumi{\arabic{enumi}.}
\item
  Enumerate the elements of length \(\leqslant 4\) in \(\mathbf{M}(\mathbf{X})\) when \(\mathbf{X}\) consists of a single element.
\item
  Let X be a set and M one of the magmas M(X), Mo(X) or \(\mathrm{N}^{(X)}\). Show that every automorphism of M leaves X stable. Deduce an isomorphism of \(\mathfrak{S}_{X}\) onto the group Aut(M).
\end{enumerate}

Show that the endomorphisms of M correspond bijectively to the mappings of X into M.

\begin{enumerate}
\def\labelenumi{\arabic{enumi}.}
\setcounter{enumi}{2}
\tightlist
\item
  Let X be a set and \(M_{a}\) the free magma constructed on a set \(\{a\}\) with one element. Let \(\rho\) denote the homomorphism of \(\mathrm{M}(\mathrm{X})\) into \(M_{a}\) which maps every element of X onto a; on the other hand let \(\lambda: \mathrm{M}(\mathrm{X}) \to \mathrm{Mo}(\mathrm{X})\) be the homomorphism defined in no. 9. Show that the mapping \(w \mapsto (\lambda(w), \rho(w))\) is an isomorphism of \(\mathrm{M}(\mathrm{X})\) onto the submagma of \(\mathrm{Mo}(\mathrm{X}) \times \mathrm{M}_{a}\) consisting of the ordered pairs \((u, v)\) where u and v have the same length.
\end{enumerate}

¶ 4. Let X be a set and c an element not belonging to X. We write Y = X ∪ \{c\}. Let M be the magma with underlying set Mo(Y) and law of composition given by (u, v) ↦ cuv. Let L be the submagma of M generated by X and let ε be the mapping from Y to Z equal to -1 on X and to 1 at c.

\begin{enumerate}
\def\labelenumi{(\alph{enumi})}
\tightlist
\item
  Show that L is the subset of M consisting of the words \(w = y_{1} \ldots y_{p}\) , \(y_{i} \in Y\) , satisfying the two following conditions:
\end{enumerate}

\[
(i) \varepsilon (y _ {1}) + \dots + \varepsilon (y _ {p}) = - 1
\]

\[
(ii) \varepsilon \left(y _ {1}\right) + \dots + \varepsilon \left(y _ {i}\right) \geqslant 0 \quad \text { for } 1 \leqslant i \leqslant p - 1.
\]

Show that every element of L - X may be written uniquely in the form cuv with \(u, v \in L\) .

(The elements \(w\) satisfying (i) and (ii) form a submagma \(L'\) of \(M\) containing \(X\) . If \(w = y_1 \ldots y_p\) belongs to \(L'\) , let \(k\) be the least integer such that \(\varepsilon(y_1) + \cdots + \varepsilon(y_k) = 0\) ; show that \(y_1 = c\) and that the elements \(u = y_2 \ldots y_k\) and \(v = y_{k+1} \ldots y_p\) belong to \(L'\) ; then \(w = cuv\) , whence by induction on \(p\) the relation \(w \in L\) . Then show that the relations \(u', v' \in L\) and \(w = cu'v'\) imply \(u' = u, v' = v\) .)

\begin{enumerate}
\def\labelenumi{(\alph{enumi})}
\setcounter{enumi}{1}
\tightlist
\item
  Show that the injection of X into L extends to an isomorphism of M(X) onto L.
\end{enumerate}

\begin{enumerate}
\def\labelenumi{\arabic{enumi}.}
\setcounter{enumi}{4}
\tightlist
\item
  Let X be a set with one element. If n is an integer \(\geqslant1\) , let \(u_{n}\) denote the number of elements in M(X) of length n.
\end{enumerate}

\begin{enumerate}
\def\labelenumi{(\alph{enumi})}
\tightlist
\item
  For every set Y, show that
\end{enumerate}

\[
\operatorname{Card} \left(\mathrm{M} _ {n} (\mathrm{Y})\right) = u _ {n}. \operatorname{Card} (\mathrm{Y}) ^ {n}.
\]

(Use Exercise 3.)

\begin{enumerate}
\def\labelenumi{(\alph{enumi})}
\setcounter{enumi}{1}
\tightlist
\item
  Establish the relation
\end{enumerate}

\[
u _ {n} = \sum_ {p = 1} ^ {n - 1} u _ {p} u _ {n - p} \quad \text { for } n \geqslant 2.
\]

\begin{enumerate}
\def\labelenumi{(\alph{enumi})}
\setcounter{enumi}{2}
\tightlist
\item
  *Let \(f(\mathbf{T})\) be the formal power series \(\sum_{n=1}^{\infty} u_n \mathbf{T}^n\) . Show that
\end{enumerate}

\[
f (\mathbf {T}) = \mathbf {T} + f (\mathbf {T}) ^ {2}.
\]

Deduce the formulae

\[
f (\mathrm{T}) = \frac {1}{2} (1 - (1 - 4 \mathrm{T}) ^ {1 / 2}), \quad u _ {n} = \frac {2 ^ {n - 1}}{n !} 1. 3. 5 \dots (2 n - 3) \quad \text { for } n \geqslant 2. *
\]

Obtain the last result using Exercise 11 of Set Theory, III, § 5.

¶ 6. Let X be a set.

\begin{enumerate}
\def\labelenumi{(\alph{enumi})}
\item
  Let N be a submagma of M(X) and Y = N - (N.N). The injection Y → N extends to a homomorphism u: M(Y) → N. Show that u is an isomorphism. (In other words, every submagma of a free magma is free.)
\item
  If \(x \in \mathbf{M}(\mathbf{X})\) , let \(\mathbf{M}_x\) denote the submagma of \(\mathbf{M}(\mathbf{X})\) generated by \(x\) ; by (a), \(\mathbf{M}_x\) is identified with the free magma constructed on \(x\) . Show that, if \(x, y \in \mathbf{M}(\mathbf{X})\) , then either \(\mathbf{M}_x \subset \mathbf{M}_y\) or \(\mathbf{M}_y \subset \mathbf{M}_x\) or \(\mathbf{M}_x \cap \mathbf{M}_y = \varnothing\) . (If \(\mathbf{M}_x \cap \mathbf{M}_y \neq \varnothing\) , let \(z\) be an element of \(\mathbf{M}_x \cap \mathbf{M}_y\) of minimum length; show that either \(z = x\) or \(z = y\) .)
\end{enumerate}

\begin{enumerate}
\def\labelenumi{\arabic{enumi}.}
\setcounter{enumi}{6}
\tightlist
\item
  Let \(\mathbf{M}_x\) be the free magma constructed on a set \(\{x\}\) with one element. (a) For all \(y \in \mathbf{M}_x\) , show that there exists a unique endomorphism \(f_y\) of \(\mathbf{M}_x\) such that \(f_y(x) = y\) ; it is an isomorphism of \(\mathbf{M}_x\) onto the submagma \(\mathbf{M}_y\) generated by \(y\) (cf.~Exercise 6).
\end{enumerate}

\begin{enumerate}
\def\labelenumi{(\alph{enumi})}
\setcounter{enumi}{1}
\item
  If \(y, z \in \mathbf{M}_x\) , we write \(y \circ z = f_y(z)\) . Show that the law of composition \((y, z) \mapsto y \circ z\) makes \(\mathbf{M}_x\) into a monoid with identity element \(x\) which is isomorphic to the monoid \(\operatorname{End}(\mathbf{M}_x)\) of endomorphisms of \(\mathbf{M}_x\) .
\item
  Let \(\mathbf{N}\) be the set of monogenous submagmas of \(\mathbf{M}_x\) distinct from \(\mathbf{M}_x\) . An element \(y\) of \(\mathbf{M}_x\) is called primitive if \(\mathbf{M}_y\) is a maximal element of \(\mathbf{N}\) . Show that \(xx, x(x(xx))\) are primitive whereas \(x, (xx)(xx)\) are not.
\item
  Let \(\mathbf{P}\) be the set of primitive elements of \(\mathbf{M}_x\) . Show that, if \(y, z \in \mathbf{P}, y \neq z\) , then \(\mathbf{M}_y \cap \mathbf{M}_z = \varnothing\) (use Exercise 6).
\item
  Let \(\operatorname{Mo}(\mathbf{P})\) be the free monoid constructed on \(\mathbf{P}\) . If \(\mathbf{M}_x\) is given the monoid structure defined in (b), the injection \(\mathbf{P} \to \mathbf{M}_x\) can be extended to a homomorphism \(p: \operatorname{Mo}(\mathbf{P}) \to \mathbf{M}_x\) . Show that \(p\) is an isomorphism.
\end{enumerate}

(If \(z \in \mathbf{M}_x\) , show by induction on \(l(z)\) that \(z \in \operatorname{Im}(p)\) . On the other hand, if \(y_1, \ldots, y_n, z_1, \ldots, z_m \in \mathbb{P}\) are such that \(p(y_1 \ldots y_n) = p(z_1 \ldots z_m)\) , then \(\mathbf{M}_{y_1} \cap \mathbf{M}_{z_1} \neq \varnothing\) , whence \(y_1 = z_1\) and \(f_{y_1}(p(y_2 \ldots y_n)) = f_{y_1}(p(z_2 \ldots z_m))\) ; whence \(p(y_2 \ldots y_n) = p(z_2 \ldots z_m)\) and it follows that \(y_1 \ldots y_n = z_1 \ldots z_m\) arguing by induction on \(\sup(n, m)\) .)

\begin{enumerate}
\def\labelenumi{\arabic{enumi}.}
\setcounter{enumi}{7}
\tightlist
\item
  Let X be a set; for every integer \(q \geqslant 0\) let \(\mathrm{Mo}(\mathbf{X})_{q}\) be the set of elements of \(\mathrm{Mo}(\mathbf{X})\) of length q and let \(\mathrm{Mo}^{(q)}(\mathbf{X})\) be the union of \(\mathrm{Mo}(\mathbf{X})_{nq}, n \in \mathbb{N}\) . The injection \(\mathrm{Mo}(\mathbf{X})_{q} \to \mathrm{Mo}^{(q)}(\mathbf{X})\) extends to a homomorphism
\end{enumerate}

\[
\operatorname{Mo} (\operatorname{Mo} (\mathrm{X}) _ {q}) \rightarrow \operatorname{Mo} ^ {(q)} (\mathrm{X});
\]

show that this homomorphism is bijective if \(q \geqslant 1\) .

\begin{enumerate}
\def\labelenumi{\arabic{enumi}.}
\setcounter{enumi}{8}
\item
  Let \(x, y \in \mathbf{X}\) with \(x \neq y\) . Show that the submonoid of \(\operatorname{Mo}(\mathbf{X})\) generated by \(\{x, xy, yx\}\) is not isomorphic to a free monoid.
\item
  Show that the group defined by the presentation
\end{enumerate}

\[
\langle x, y; x y ^ {2} = y ^ {3} x, y x ^ {2} = x ^ {3} y \rangle
\]

reduces to the identity element.

(The first relation implies \(x^{2}y^{8}x^{-2} = y^{18}\) and \(x^{3}y^{8}x^{-3} = y^{27}\) ; use the second relation to deduce that \(y^{18} = y^{27}\) , whence \(y^{9} = e\) ; the fact that \(y^{2}\) is conjugate to \(y^{3}\) then implies \(y = e\) , whence \(x = e\) .)

\begin{enumerate}
\def\labelenumi{\arabic{enumi}.}
\setcounter{enumi}{10}
\tightlist
\item
  The group defined by the presentation \(\langle x,y;x^{2},y^{3},(xy)^{2}\rangle\) is isomorphic to \(S_{3}\) .
\end{enumerate}

¶ 12. The group defined by the presentation \(\langle x, y; x^3, y^2, (xy)^3 \rangle\) is isomorphic to \(\mathfrak{A}_4\) . (Show that the conjugates of \(y\) form, with the identity element, a normal subgroup of order \(\leqslant 4\) and index \(\leqslant 3\) .)

\begin{enumerate}
\def\labelenumi{\arabic{enumi}.}
\setcounter{enumi}{12}
\tightlist
\item
  Let \(G\) be a group and \(X\) a subset of \(G\) disjoint from \(X^{-1}\) . Let \(S = X \cup X^{-1}\) . Show the equivalence of the two following properties:
\end{enumerate}

\begin{enumerate}
\def\labelenumi{(\roman{enumi})}
\item
  X is a free family in G.
\item
  No product \(s_1 \ldots s_n\) , with \(n \geqslant 1\) , \(s_i \in S\) and \(s_i s_{i+1} \neq e\) for \(1 \leqslant i \leqslant n-1\) , is equal to \(e\) .
\end{enumerate}

\begin{enumerate}
\def\labelenumi{\arabic{enumi}.}
\setcounter{enumi}{13}
\tightlist
\item
  A group G is called free if it possesses a basic family and hence is isomorphic to a free group F(X) constructed on a set X.
\end{enumerate}

\begin{enumerate}
\def\labelenumi{(\alph{enumi})}
\tightlist
\item
  Show that, if \((t_i)_{i\in I}\) and \((t_j')_{j\in J}\) are two such families, then \(\operatorname{Card}(I) = \operatorname{Card}(J)\) . (When \(\operatorname{Card}(I)\) is infinite, show that \(\operatorname{Card}(J)\) is also and that both are equal to \(\operatorname{Card}(G)\) . When \(\operatorname{Card}(I)\) is an integer \(d\) , show that the number of subgroups of index 2 in \(G\) is \(2^d - 1\) .)
\end{enumerate}

The cardinal of a basic family of G is called the rank of G.

\begin{enumerate}
\def\labelenumi{(\alph{enumi})}
\setcounter{enumi}{1}
\item
  A free group is of rank 0 (resp. 1) if and only if it reduces to \(\{e\}\) (resp. is isomorphic to \(\mathbf{Z}\) ).
\item
  Let G be a free group of finite rank d and \((x_{1},\ldots,x_{d})\) a generating family of d elements of G. Show that this family is basic. (Use Exercise 34 and § 5, Exercise 5.)
\item
  Show that a free group of rank \(\geqslant 2\) contains a free subgroup of given rank \(d\) for every cardinal \(d \leqslant \aleph_0\) . (Use Exercise 22.)
\end{enumerate}

¶ 15. Let G be a group with presentation (t, r), where \(\mathbf{t} = (t_{i})_{i \in I}\) and \(\mathbf{r} = (r_{j})_{j \in J}\) ; let \(\mathrm{F}((\mathrm{T}_{i})_{i \in I})\) denote the free group \(\mathrm{F}(I)\) and \(\mathbf{T} = (\mathrm{T}_{i})\) . On the other hand, let \(x = (x_{\alpha})_{\alpha \in A}\) be a generating family of G; write

\[
\mathrm{F} (\mathbf {A}) = \mathrm{F} ((\mathbf {X} _ {\alpha}) _ {\alpha \in A}),
\]

\(\mathbf{X} = (\mathbf{X}_{\alpha})\) and \(\mathbf{R}_{\mathbf{x}}\) the set of relators of \(\mathbf{x}\) ; it is a normal subgroup of \(\mathrm{F(A)}\) . For all \(i \in \mathbf{I}\) , let \(\phi_i(\mathbf{X})\) be an element of \(\mathrm{F(A)}\) such that \(t_i = \phi_i(\mathbf{x})\) in \(\mathbf{G}\) ; for all \(\alpha \in \mathbf{A}\) , let \(\psi_{\alpha}(\mathbf{T})\) be an element of \(\mathrm{F(I)}\) such that \(x_{\alpha} = \psi_{\alpha}(\mathbf{t})\) in \(\mathbf{G}\) . Show that \(\mathbf{R}_{\mathbf{x}}\) is the normal subgroup of \(\mathrm{F(A)}\) generated by the elements \(\mathbf{X}_{\alpha}^{-1} \psi_{\alpha}(\phi_i(\mathbf{X})_{i \in I}), \alpha \in \mathbf{A}\) and \(r_j(\phi_i(\mathbf{X})_{i \in I}) j \in J\) . (If \(\mathbf{R}_{\mathbf{x}}'\) denotes the normal subgroup of \(\mathrm{F(A)}\) generated by the elements in question, then \(\mathbf{R}_{\mathbf{x}}' \subset \mathbf{R}_{\mathbf{x}}\) ; on the other hand, show that the homomorphism \(\mathrm{F(I)} \to \mathrm{F(A)}\) defined by the \(\phi_i\) gives, when passing to the quotient, a homomorphism \(\mathrm{G} \to \mathrm{F(A)} / \mathrm{R}_{\mathbf{x}}'\) the inverse of the canonical homomorphism \(\mathrm{F(A)} / \mathrm{R}_{\mathbf{x}}' \to \mathrm{F(A)} / \mathrm{R}_{\mathbf{x}} = \mathrm{G}\) .)

\begin{enumerate}
\def\labelenumi{\arabic{enumi}.}
\setcounter{enumi}{15}
\tightlist
\item
  A group is called finitely generated (resp. finitely presented) if it admits a presentation \(\langle (t_i)_{i\in I}, (r_j)_{j\in J}\rangle\) where \(I\) is a finite set (resp. where \(I\) and \(J\) are finite sets).
\end{enumerate}

\begin{enumerate}
\def\labelenumi{(\alph{enumi})}
\item
  Let G be a finitely presented group and \(\mathbf{x} = (x_{\alpha})_{\alpha \in \mathbf{A}}\) a finite generating family of G. Show that there exists a finite family \(\mathbf{s} = (s_{k})_{k \in \mathbf{K}}\) of relators of the family x such that \((\mathbf{x}, \mathbf{s})\) is a presentation of G. (Use the preceding exercise.)
\item
  Show that a finite group is finitely presented.
\item
  Give an example of a finitely presented group with a subgroup which is not finitely generated.
\item
  If \(\mathbf{G}\) is a finitely presented group and \(\mathbf{H}\) a normal subgroup of \(\mathbf{G}\) , show the equivalence of the following properties:
\item
  G/H is finitely presented.
\end{enumerate}

\begin{enumerate}
\def\labelenumi{(\roman{enumi})}
\setcounter{enumi}{1}
\item
  There exists a finite subset X of G such that H is the normal subgroup of G generated by X.
\item
  If \((\mathbf{H}_i)_{i\in \mathbf{I}}\) is a right directed family of normal subgroup of \(\mathbf{G}\) whose union is equal to \(\mathbf{H}\) , there exists \(i\in \mathbf{I}\) such that \(\mathbf{H}_i = \mathbf{H}\) .
\end{enumerate}

(Use (a) to prove that (i) implies (ii).)

\begin{enumerate}
\def\labelenumi{(\alph{enumi})}
\setcounter{enumi}{4}
\item
  If G is a finitely presented group, show that \(\mathbf{G}/\mathbf{C}^{r}(\mathbf{G})\) is finitely presented for all \(r \geqslant 1\) . (Use §6, Exercise 15.)
\item
  Let \((\mathbf{G}_{\alpha})_{\alpha \in \mathbf{A}}\) be a family of groups and let \(\mathbf{G} = \prod_{\alpha \in \mathbf{A}} \mathbf{G}_{\alpha}\) be the product of the \(\mathbf{G}_{\alpha}\) . Show that \(\mathbf{G}\) is finitely generated (resp. finitely presented) if and only if each of the \(\mathbf{G}_{\alpha}\) is finitely generated (resp. finitely presented) and \(\mathbf{G}_{\alpha} = \{e\}\) for all but a finite number of \(\alpha\) .
\end{enumerate}

¶ 17. (a) Show that every subgroup of a finitely generated nilpotent group is finitely generated. (In the commutative case, argue by induction on the minimum number of generators of the group; then proceed by induction on the nilpotency class of the group.)

Give an example of a finitely generated solvable group containing a subgroup which is not finitely generated.

\begin{enumerate}
\def\labelenumi{(\alph{enumi})}
\setcounter{enumi}{1}
\tightlist
\item
  Show that every finitely generated nilpotent group is finitely presented. (Write the group in the form \(\mathrm{F}(\mathbf{X}) / \mathbb{R}\) , with \(\mathbf{X}\) finite, and choose \(r\) such that \(\mathbf{R} \supset \mathbf{C}^r (\mathrm{F}(\mathbf{X}))\) ; use (a) to show that \(\mathbb{R} / \mathbb{C}^r (\mathrm{F}(\mathbf{X}))\) is finitely generated; observe that \(\mathrm{F}(\mathbf{X}) / \mathbb{C}^r (\mathrm{F}(\mathbf{X}))\) is finitely presented, cf.~Exercise 16.)
\end{enumerate}

\begin{enumerate}
\def\labelenumi{\arabic{enumi}.}
\setcounter{enumi}{17}
\tightlist
\item
  Let S be a symmetric subset of a group G. Two elements \(g_{1}, g_{2}\) of G are called S-neighbours if \(g_{2}^{-1}g_{1}\) belongs to S. A sequence \((g_{1}, \ldots, g_{n})\) of elements of G is called an S-chain if \(g_{i}\) and \(g_{i+1}\) are S-neighbours for \(1 \leqslant i \leqslant n - 1\) ; the elements \(g_{1}\) and \(g_{n}\) are called respectively the beginning and the end of the chain.
\end{enumerate}

\begin{enumerate}
\def\labelenumi{(\alph{enumi})}
\item
  Let Y be a subset of G and \(R_{Y}\{a, b\}\) the relation ``there exists an S-chain in Y beginning at a and ending at \(b''\) . Show that \(R_{Y}\) is an equivalence relation on Y (cf.~Set Theory, II, §6, Exercise 10); the equivalence classes under this relation are called the connected S-components of Y. Y is called S-connected if it has at most one connected component.
\item
  Show that \(\mathbf{G}\) is S-connected if and only if \(\mathbf{S}\) generates \(\mathbf{G}\) .
\item
  Suppose that no element \(s \in S\) satisfies the relation \(s^2 = e\) ; choose a subset \(X\) of \(S\) such that \(S\) is the disjoint union of \(X\) and \(X^{-1}\) . Show the equivalence of the following conditions:
\item
  The family \(\mathbf{X}\) is free.
\end{enumerate}

\begin{enumerate}
\def\labelenumi{(\roman{enumi})}
\setcounter{enumi}{1}
\tightlist
\item
  There exists no S-chain \((g_{1},\ldots ,g_{n})\) in G consisting of distinct elements \(n\geqslant 3\) in number such that \(g_{1}\) and \(g_{n}\) are S-neighbours.
\end{enumerate}

(Use Exercise 13.)

¶ 19. Let X be a set, G = F(X) the free group constructed on X and S = X ∪ X \(^{-1}\) .

\begin{enumerate}
\def\labelenumi{(\alph{enumi})}
\item
  If \(g \in G\) , every sequence \((s_1, \ldots, s_n)\) of elements of \(S\) such that \(g = s_1 \ldots s_n\) and \(s_i s_{i+1} \neq e\) for \(1 \leqslant i < n\) is called a reduced decomposition of \(g\) . Show that every element of \(G\) admits one and only one reduced decomposition; the integer \(n\) is the length of \(g\) (cf.~Exercise 26).
\item
  Let Y be a subset of G containing e. Show the equivalence of the following conditions:
\end{enumerate}

\((b_{1})\) Y is S-connected (cf.~Exercise 18).

\((b_{2})\) For all \(g \in \mathbf{Y}\) , if \((s_1, \ldots, s_n)\) is the reduced decomposition of \(g\) , then \(s_1 \ldots s_i \in \mathbf{Y}\) for \(1 \leqslant i \leqslant n\) .

\begin{enumerate}
\def\labelenumi{(\alph{enumi})}
\setcounter{enumi}{2}
\tightlist
\item
  Let \(Y_{1}\) and \(Y_{2}\) be two non-empty S-connected subsets of G such that \(Y_{1} \cap Y_{2} = \varnothing\) . Show that there exists at most one ordered pair
\end{enumerate}

\[
\left(y _ {1}, y _ {2}\right) \in \mathbf {Y} _ {1} \times \mathbf {Y} _ {2}
\]

such that \(y_{1}\) and \(y_{2}\) are S-neighbours. There exists one if and only if \(\mathbf{Y}_1 \cup \mathbf{Y}_2\) is S-connected; \(\mathbf{Y}_1\) and \(\mathbf{Y}_2\) are then said to be neighbours.

\begin{enumerate}
\def\labelenumi{(\alph{enumi})}
\setcounter{enumi}{3}
\tightlist
\item
  Let \(Y_{1}, \ldots, Y_{n}\) be a sequence of disjoint non-empty S-connected subsets of G. Suppose that \(Y_{i}\) and \(Y_{i+1}\) are neighbours for \(1 \leqslant i < n\) . Show that, if \(n \geqslant 3\) , \(Y_{1}\) and \(Y_{n}\) are not neighbours.
\end{enumerate}

¶ 20. We preserve the notation of the preceding exercise. Let H be a subgroup of G.

\begin{enumerate}
\def\labelenumi{(\alph{enumi})}
\item
  Let \(R_{H}\) be the set of S-connected subsets T of G such that \(hT \cap T = \varnothing\) if \(h \in H, h \neq e\) . Show that \(R_{H}\) is inductive for the relation of inclusion.
\item
  Let Y be a maximal element of \(R_{H}\) . Show that \(G = \bigcup_{h \in H} hY\) , in other words Y is a system of representatives of the right cosets modulo H. (If \(G' = \bigcup hY\) is distinct from G, show that there exist \(x \in G'\) and \(y \in Y\) which are S-neighbours and deduce that \(Y \cup \{x\}\) belongs to \(R_{H}\) .)
\item
  Let \(S_{H}\) be the set of elements \(h \in H\) such that Y and hY are neighbours. It is a symmetric subset of H. Show that an element \(h \in H\) belongs to \(S_{H}\) if and only if \(h \neq e\) and there exist \(y, y' \in Y, s \in S\) with \(ys = hy'\) .
\end{enumerate}

Show that \(S_{H}\) generates H. (If \(H'\) is the subgroup of H generated by \(S_{H}\) .

prove that \(\bigcup_{h\in H'}hY\) is a connected S-component of G.) Obtain this result by means of Exercise 14 of §4.

\begin{enumerate}
\def\labelenumi{(\alph{enumi})}
\setcounter{enumi}{3}
\tightlist
\item
  Let \(X_{H}\) be a subset of \(S_{H}\) such that \(S_{H}\) is the disjoint union of \(X_{H}\) and \(X_{H}^{-1}\) (show that such a subset exists).
\end{enumerate}

Show that \(X_{H}\) is a basic family for H and in particular that H is a free group (``Nielsen-Schreier Theorem''). (Apply the criterion of Exercise 18(c) to H together with \(S_{H}\) , noting that two elements h, \(h'\) of H are \(S_{H}\) -neighbours if and only if hY and \(h'Y\) are neighbouring subsets of G. Use Exercise 19(d) to prove that \(S_{H}\) satisfies condition (ii) in Exercise 18(c).)

\begin{enumerate}
\def\labelenumi{(\alph{enumi})}
\setcounter{enumi}{4}
\tightlist
\item
  Let \(d = (\mathbf{G}:\mathbf{H}) = \operatorname{Card}(\mathbf{Y})\) . Suppose \(d\) is finite. Show that:
\end{enumerate}

\[
\begin{array}{l l} \operatorname{Card} (\mathbf {X} _ {\mathrm{H}}) = \operatorname{Card} (\mathbf {X}) = \operatorname{Card} (\mathbf {G}) & \text { if } \operatorname{Card} (\mathbf {X}) \text { is   infinite } \\ \operatorname{Card} (\mathbf {X} _ {\mathrm{H}}) = 1 + d (\operatorname{Card} (\mathbf {X}) - 1) & \text { if } \operatorname{Card} (\mathbf {X}) \text { is   finite }. \end{array}
\]

(Let T be a non-empty S-connected finite subset of G and T' (resp. T'') the set of ordered pairs of S-neighbouring elements whose first element (resp. both of whose elements) belongs (resp. belong) to T. Card(T') = 2Card(X)Card(T) and it should be shown by induction on Card(T) that

\[
\operatorname{Card} \left(\mathrm{T} ^ {\prime \prime}\right) = 2 \operatorname{Card} (\mathrm{T}) - 2.
\]

Apply the result to T = Y, noting that \(S_{H}\) is equipotent to \(T' - T''\) .

\begin{enumerate}
\def\labelenumi{\arabic{enumi}.}
\setcounter{enumi}{20}
\item
  Let H be a subgroup of a group G of finite index. Show that H is finitely presented if and only if G is. (It may be assumed that G is finitely generated, cf.~§4, Exercise 14. Write G in the form G = F/R with F free and finitely generated; the inverse image F' of H in F is free and finitely generated, cf.~Exercise 20, and H = F'/R. If R is generated (as a normal subgroup of F) by \((r_{j})_{j\in J}\) , show that it is generated (as a normal subgroup of F') by the \(yr_{j}y^{-1}\) , \(y\in Y\) , \(j\in J\) , where Y denotes a representative system in F of the right cosets modulo F'.)
\item
  \begin{enumerate}
  \def\labelenumii{(\alph{enumii})}
  \tightlist
  \item
    Let \((y_{n})_{n\in \mathbf{Z}}\) be a basic family of a group F. Let \(\phi\) be the automorphism of F which maps \(y_{n}\) to \(y_{n + 1}\) ; this automorphism defines an operation \(\tau\) of \(\mathbf{Z}\) on F; let \(\mathbf{E} = \mathbf{F}\times_{\tau}\mathbf{Z}\) be the corresponding semi-direct product. Show that the elements \(x = (e,1)\) and \(y = (y_0,0)\) form a basic family of E. (Let \(\mathrm{F}_{x,y}\) be the free group constructed on \(\{x,y\}\) and let \(f\) be the canonical homomorphism of \(\mathrm{F}_{x,y}\) into E. Show that there exists a homomorphism \(g\colon \mathbf{E}\to \mathbf{F}_{x,y}\) such that \(g((0,1)) = x\) and \(g((y_n,0)) = x^n yx^{-n}\) and \(f\) and \(g\) are inverses of one another.)
  \end{enumerate}
\end{enumerate}

\begin{enumerate}
\def\labelenumi{(\alph{enumi})}
\setcounter{enumi}{1}
\item
  Deduce that the normal subgroup of \(F_{x,y}\) generated by y is a free group with basic family \((x^{n}yx^{-n})_{n\in\mathbf{Z}}\) . Obtain this result by applying Exercise 20 to the representative system consisting of the \(x^{n}, n\in Z\) .
\item
  Extend the above to free groups of arbitrary rank.
\end{enumerate}

\begin{enumerate}
\def\labelenumi{\arabic{enumi}.}
\setcounter{enumi}{22}
\tightlist
\item
  Let \(\mathbf{F}\) be a free group with basic family \((x, y)\) , \(x \neq y\) . Show that the derived group of F has as basic family the family of commutators \((x^i, y^j)\) , \(i \in \mathbf{Z}\) , \(j \in \mathbf{Z}\) , \(i \neq 0\) , \(j \neq 0\) .
\end{enumerate}

(Use Exercise 20 or Exercise 32.)

\begin{enumerate}
\def\labelenumi{\arabic{enumi}.}
\setcounter{enumi}{23}
\tightlist
\item
  Let G be a group, \(S = G - \{e\}\) and \(F = F(S)\) . For all \(s \in S\) , let \(x_s\) denote the corresponding element of F. If \(s, t \in S\) , let \(r_{s,t}\) denote the element of F defined by:
\end{enumerate}

\[
r _ {s, t} = x _ {s} x _ {t} x _ {s t} ^ {- 1} \quad \text { if } s t \neq e \text { in } G
\]

\[
r _ {s, t} = x _ {s} x _ {t} \quad \text { if } s t = e \text { in } G.
\]

Let \(\phi\) be the homomorphism of F onto G which maps \(x_{s}\) to \(s\) and let R be the kernel of \(\phi\) .

Show that the family \((r_{s,t})\) for \(s, t \in \mathbf{S} \times \mathbf{S}\) , is a basic family of the group \(\mathbf{R}\) . (Apply Exercise 20 to the system of representatives of \(\mathbf{G}\) in \(\mathbf{F}\) consisting of \(e\) and the \(x_{s'}\) .)

\begin{enumerate}
\def\labelenumi{\arabic{enumi}.}
\setcounter{enumi}{24}
\tightlist
\item
  Let G be a group. Show the equivalence of the following properties:
\end{enumerate}

\begin{enumerate}
\def\labelenumi{(\alph{enumi})}
\item
  G is a free group.
\item
  For every group H and extension E of G by H there exists a section \(G \rightarrow E\) .
\end{enumerate}

(To prove that (b) implies (a), take E to be a free group and use Exercise 20.)

\begin{enumerate}
\def\labelenumi{\arabic{enumi}.}
\setcounter{enumi}{25}
\tightlist
\item
  Let \(\mathbf{X}\) be a set and \(g\) an element of the free group \(\mathrm{F}(\mathbf{X})\) . Let \(g = \prod_{\alpha=1}^{n} x_{\alpha}^{m(\alpha)}\) be the canonical decomposition of \(g\) as a product of powers of elements of \(\mathbf{X}\) (cf.~no. 5, Proposition 7). The integer \(l(g) = \sum_{\alpha} |m(\alpha)|\) is called the length of \(g\) .
\end{enumerate}

\begin{enumerate}
\def\labelenumi{(\alph{enumi})}
\item
  \(g\) is called cyclically reduced if either \(x_{1} \neq x_{n}\) or \(x_{1} = x_{n}\) and \(m(1)m(n) > 0\) . Show that every conjugacy class of \(F(X)\) contains one and only one cyclically reduced element; it is the element of minimum length of the class in question.
\item
  An element \(x \in \mathrm{F}(\mathbf{X})\) is called primitive if there does not exist an integer \(n \geqslant 2\) and an element \(g \in \mathrm{F}(\mathbf{X})\) such that \(x = g^n\) . Show that every element \(z \neq e\) of \(\mathrm{F}(\mathbf{X})\) can be written uniquely in the form \(x^n\) with \(n \geqslant 1\) and \(x\) primitive. (Reduce it to the case where \(z\) is cyclically reduced and observe that, if \(z = x^n\) , the element \(x\) is also cyclically reduced.)
\end{enumerate}

Show that the centralizer of z is the same as that of x; it is the cyclic subgroup generated by x.

\begin{enumerate}
\def\labelenumi{\arabic{enumi}.}
\setcounter{enumi}{26}
\tightlist
\item
  Let \(G = \times_{A} G_{i}\) be the sum of a family \((G_{i})_{i \in I}\) of groups amalgamated by a common subgroup A. For all \(i \in I\) , choose a subset \(P_{i}\) of \(G_{i}\) satisfying condition (A) of no. 3.
\end{enumerate}

\begin{enumerate}
\def\labelenumi{(\alph{enumi})}
\tightlist
\item
  Let \(x \in G\) and let \((a; i_1, \ldots, i_n; p_1, \ldots, p_n)\) be a reduced decomposition of \(x\) ; then
\end{enumerate}

\[
x = a \prod_ {\alpha = 1} ^ {n} p _ {\alpha}, \quad \text { with } \quad p _ {\alpha} \in \mathrm{P} _ {i _ {\alpha}} - \{e _ {i _ {\alpha}} \}, i _ {\alpha} \neq i _ {\alpha + 1}.
\]

Show that, if \(i_1 \neq i_n\) , the order of \(x\) is infinite. Show that, if \(i_1 = i_n\) , \(x\) is conjugate to an element with a decomposition of length \(\leqslant n - 1\) .

\begin{enumerate}
\def\labelenumi{(\alph{enumi})}
\setcounter{enumi}{1}
\tightlist
\item
  Deduce that every element of G of finite order is conjugate to an element of one of the \(G_{i}\) . In particular, if the \(G_{i}\) have no element of finite order except e, the same is true of G.
\end{enumerate}

\begin{enumerate}
\def\labelenumi{\arabic{enumi}.}
\setcounter{enumi}{27}
\tightlist
\item
  We preserve the notation of the preceding exercise.
\end{enumerate}

\begin{enumerate}
\def\labelenumi{(\alph{enumi})}
\item
  Let \(\mathbf{N}\) be a subgroup of \(\mathbf{A}\) . Suppose that, for all \(i \in \mathbf{I}\) , the group \(\mathbf{N}\) is normal in \(\mathbf{G}_i\) . Show that \(\mathbf{N}\) is normal in \(\mathbf{G} = \times_{\mathbf{A}} \mathbf{G}_i\) and that the canonical homomorphism of \(\times_{\mathbf{A/N}} \mathbf{G}_i / \mathbf{N}\) into \(\mathbf{G}/\mathbf{N}\) is an isomorphism.
\item
  For all \(i \in \mathbf{I}\) , let \(\mathbf{H}_i\) be a subgroup of \(\mathbf{G}_i\) and let \(\mathbf{B}\) be a subgroup of \(\mathbf{A}\) . Suppose that \(\mathbf{H}_i \cap \mathbf{A} = \mathbf{B}\) for all \(i\) . Show that the canonical homomorphism of \(\star_B \mathbf{H}_i\) into \(\mathbf{G}\) is injective; its image is the subgroup of \(\mathbf{G}\) generated by the \(\mathbf{H}_i\) .
\end{enumerate}

¶ 29. Let \(G = G_1 *_A G_2\) be the sum of two groups \(G_1\) and \(G_2\) amalgamated by a subgroup A.

\begin{enumerate}
\def\labelenumi{(\alph{enumi})}
\item
  Show that G is finitely generated if \(G_{1}\) and \(G_{2}\) are finitely generated.
\item
  Suppose \(\mathbf{G}_1\) and \(\mathbf{G}_2\) are finitely presented. Show the equivalence of the two following properties:
\end{enumerate}

\begin{enumerate}
\def\labelenumi{(\arabic{enumi})}
\item
  G is finitely presented.
\item
  A is finitely generated.
\end{enumerate}

(Show first that (1) implies (2). On the other hand, if A is not finitely generated, there exists a right directed family \((\mathbf{A}_{i})_{i\in I}\) of subgroups of A with \(\bigcup_{i\in I}\mathbf{A}_{i}=\mathbf{A}\) and \(A_{i}\neq A\) for all i. Let \(H_{i}\) (resp. H) be the kernel of the canonical homomorphism of \(G_{1}*G_{2}\) onto \(G_{1}*_{A_{i}}G_{2}\) (resp. onto G). Show that H is the union of the \(H_{i}\) and that \(H_{i}\neq H\) for all i; then apply Exercise 16(d).

\begin{enumerate}
\def\labelenumi{(\alph{enumi})}
\setcounter{enumi}{2}
\tightlist
\item
  Deduce an example of a finitely generated group which is not finitely presented. (Take \(G_{1}\) and \(G_{2}\) to be free groups of rank 2 and A a free group of finite rank, cf.~Exercise 22.)
\end{enumerate}

\begin{enumerate}
\def\labelenumi{\arabic{enumi}.}
\setcounter{enumi}{29}
\tightlist
\item
  Let A and B be two groups and G = A * B their free product.
\end{enumerate}

\begin{enumerate}
\def\labelenumi{(\alph{enumi})}
\tightlist
\item
  Let \(\mathbf{Z}\) be an element of \(\mathbf{A} \backslash \mathbf{G} / \mathbf{A}\) distinct from \(\mathbf{A}\) . Show that \(\mathbf{Z}\) contains one and only one element \(z\) of the form
\end{enumerate}

\[
b _ {1} a _ {1} b _ {2} a _ {2} \dots b _ {n - 1} a _ {n - 1} b _ {n},
\]

with \(n \geqslant 1\) , \(a_i \in \mathbf{A} - \{e\}\) , \(b_i \in \mathbf{B} - \{e\}\) . Show that every element of \(\mathbf{Z}\) can be written uniquely in the form \(aza'\) with \(a, a' \in \mathbf{A}\) .

\begin{enumerate}
\def\labelenumi{(\alph{enumi})}
\setcounter{enumi}{1}
\tightlist
\item
  Let \(x \in \mathbf{G} - \mathbf{A}\) . Let \(f_x\) be the homomorphism of \(\mathbf{A} * \mathbf{A}\) into \(\mathbf{G}\) whose restriction to the first (resp. second) factor of \(\mathbf{A} * \mathbf{A}\) is the identity (resp. the mapping \(a \mapsto xax^{-1}\) ). Show that \(f_x\) is injective. (Write \(x\) in the form \(aza'\) as above and reduce it thus to the case \(x = z\) ; use the uniqueness of reduced decompositions in \(\mathbf{G}\) .)
\end{enumerate}

In particular, \(\mathbf{A} \cap x\mathbf{A}x^{-1} = \{e\}\) and the normalizer of \(\mathbf{A}\) in \(\mathbf{G}\) is \(\mathbf{A}\) .

\begin{enumerate}
\def\labelenumi{(\alph{enumi})}
\setcounter{enumi}{2}
\tightlist
\item
  Let \(\mathbf{T}\) be an element of \(\mathbf{A} \backslash \mathbf{G} / \mathbf{B}\) . Show that \(\mathbf{T}\) contains one and only one element \(t\) of the form
\end{enumerate}

\[
b _ {1} a _ {1} b _ {2} a _ {2} \dots b _ {n - 1} a _ {n - 1},
\]

with \(n \geqslant 1\) , \(a_i \in \mathbf{A} - \{e\}\) , \(b_i \in \mathbf{B} - \{e\}\) . Show that every element of \(\mathbf{T}\) can be written uniquely in the form atb with \(a \in \mathbf{A}\) , \(b \in \mathbf{B}\) .

\begin{enumerate}
\def\labelenumi{(\alph{enumi})}
\setcounter{enumi}{3}
\tightlist
\item
  Let \(x \in G\) and let \(h_x\) be the unique endomorphism of \(G\) such that \(h_x(a) = a\) if \(a \in A\) and \(h_x(b) = x b x^{-1}\) if \(b \in B\) . Show that \(h_x\) is injective. (Same method as in (b), using the decomposition \(x = a t b\) of (c).) In particular, \(A \cap x B x^{-1} = \{e\}\) .
\end{enumerate}

Give an example where \(h_{x}\) is not surjective.

\begin{enumerate}
\def\labelenumi{\arabic{enumi}.}
\setcounter{enumi}{30}
\tightlist
\item
  Let G be the free product of a family \((G_{i})_{i\in I}\) of groups. Let S be the set of subgroups of G which are conjugate to one of the \(G_{i}\) and \(\Theta\) the union of the elements of S.
\end{enumerate}

\begin{enumerate}
\def\labelenumi{(\alph{enumi})}
\item
  Let \(\mathbf{H}, \mathbf{H}' \in \mathbf{S}\) , with \(\mathbf{H} \neq \mathbf{H}'\) , and let \(x \in \mathbf{H} - \{e\}\) , \(x' \in \mathbf{H}' - \{e\}\) . Show that \(xx' \in \mathbf{G} - \Theta\) . (Reduce it to the case where \(\mathbf{H}\) is one of the \(\mathbf{G}_i\) ; write \(\mathbf{H}'\) in the form \(y\mathrm{G}_j y^{-1}\) and proceed as in the preceding exercise.)
\item
  Let \(\mathbf{C}\) be a subgroup of \(\mathbf{G}\) contained in \(\Theta\) . Show that there exists \(\mathbf{H} \in \mathbf{S}\) such that \(\mathbf{C} \subset \mathbf{H}\) .
\item
  Let C be a subgroup of G all of whose elements are of finite order. Show that C is contained in \(\Theta\) (cf.~Exercise 27); deduce that C is contained in a conjugate of one of the \(G_{i}\) .
\end{enumerate}

¶ 32. Let A and B be two groups and G = A * B their free product. Let X be the set of commutators \((a, b)\) with \(a \in A - \{e\}\) and \(b \in B - \{e\}\) and \(S = X \cup X^{-1}\) .

\begin{enumerate}
\def\labelenumi{(\alph{enumi})}
\item
  Show that \(\mathbf{X} \cap \mathbf{X}^{-1} = \varnothing\) and that, if \(x\) is an element of \(\mathbf{X}\) , there exists a unique ordered pair \(a, b\) in \((\mathbf{A} - \{e\}) \times (\mathbf{B} - \{e\})\) such that \((a, b) = x\) .
\item
  Let \(n \geqslant 1\) and let \(s_1, \ldots, s_n\) be a sequence of elements of \(S\) such that \(s_i s_{i+1} \neq e\) for \(1 \leqslant i < n\) ; let \(a_n \in A - \{e\}, b_n \in B - \{e\}, \varepsilon(n) = \pm 1\) be such that \(s_n = (a_n, b_n)^{\varepsilon(n)}\) . Let \(g = s_1 \ldots s_n\) . Show, by induction on \(n\) , that the reduced decomposition of \(g\) is of length \(\geqslant n + 3\) and terminates either with \(a_nb_n\) if \(\varepsilon(n) = 1\) or with \(b_na_n\) if \(\varepsilon(n) = -1\) .
\item
  Let R be the kernel of the canonical projection \(A*B \rightarrow A \times B\) corresponding (no. 3) to the canonical injections \(A \rightarrow A \times B\) and \(B \rightarrow A \times B\) . Show that R is a free group with basic family X. (Use (b) to prove that X is free; if \(R_{X}\) is the subgroup of G generated by X, show that \(R_{X}\) is normal; deduce that it coincides with R.)
\end{enumerate}

\begin{enumerate}
\def\labelenumi{\arabic{enumi}.}
\setcounter{enumi}{32}
\tightlist
\item
  Let \(E = G *_{A} G'\) be the sum of two groups amalgamated by a common subgroup A. Let P (resp. \(P'\) ) be a subset of G (resp. \(G'\) ) satisfying condition (A) of no. 3.
\end{enumerate}

Let \(n\) be an integer \(\geqslant 1\) . Let \(L_n\) denote the set of \(x \in E\) whose reduced decomposition is of length \(\leqslant 2n - 1\) . Let \(\mathbf{M}_n\) denote the set of \(x \in \mathbf{E}\) of the form \(ap_1p_1'p_2p_2'\ldots p_np_n'\) , with \(a \in \mathbf{A}, p_i \in \mathbf{P} - \{e\}, p_i' \in \mathbf{P}' - \{e\}\) ; similarly, let \(\mathbf{M}_n'\) denote the set of elements of the form \(ap_1'p_1p_2'p_2\ldots p_np_n'\) , where \(a, p_i, p_i'\) satisfy the same conditions.

\begin{enumerate}
\def\labelenumi{(\alph{enumi})}
\item
  Show that \(L_{n}\) , \(M_{n}\) and \(M_{n}^{\prime}\) are disjoint and that their union is the set of elements of E whose reduced decomposition is of length \(\leqslant2n\) .
\item
  Construct a bijection \(\varepsilon\) of \(\mathbf{M}_n\) onto \(\mathbf{M}_n'\) such that \(\varepsilon(ax) = a\varepsilon(x)\) if \(a \in \mathbf{A}\) , \(x \in \mathbf{M}_n\) .
\item
  Let \(\mathbf{X}_n = \mathbf{L}_n \cup \mathbf{M}_n\) and \(\mathbf{X}_n' = \mathbf{L}_n \cup \mathbf{M}_n'\) ; show that \(\mathbf{G} \cdot \mathbf{X}_n = \mathbf{X}_n\) and \(\mathbf{G}' \cdot \mathbf{X}_n' = \mathbf{X}_n'\) .
\item
  \(\varepsilon\) extends to a bijection of \(\mathbf{X}_n\) onto \(\mathbf{X}_n'\) by setting \(\varepsilon(x) = x\) if \(x \in \mathbf{L}_n\) . Show that E can operate on \(\mathbf{X}_n\) so that \(g(x) = gx\) if \(g \in G\) , \(x \in \mathbf{X}_n\) and \(g'(x) = \varepsilon^{-1}(g' \varepsilon (x))\) if \(g' \in G'\) , \(x \in \mathbf{X}_n\) . Show that, if \(g \in \mathbf{X}_n\) , then \(g(e) = g\) .
\item
  Let \(\tau_{n}\colon\mathbf{E}\to\mathfrak{S}_{\mathbf{X}_{n}}\) be the homomorphism defined by the action of \(\mathbf{E}\) on \(\mathbf{X}_{n}\) described above and let \(\mathbf{R}_{n}\) be the kernel of \(\tau_{n}\) . Show that \(\mathbf{R}_{n}\cap\mathbf{X}_{n}=\{e\}\) ; deduce that the intersection of the \(\mathbf{R}_{n}\) reduces to \(e\) .
\item
  When G and G' are finite, show that E is residually finite (cf.~§ 5, Exercise 5). (Note that the R \(_{n}\) are then subgroups of finite index.)
\end{enumerate}

¶ 34. (a) Let \((\mathrm{H}_{i})\) be a right directed family of normal subgroups of a group G. Suppose that \(\bigcap_{i} \mathrm{H}_{i} = \{e\}\) . Show that, for every group \(\mathbf{G}'\) , the intersection of the kernels of the canonical homomorphisms \(\mathbf{G}' * \mathbf{G} \to \mathbf{G}' * (\mathbf{G}/\mathbf{H}_{i})\) is equal to \(\{e\}\) .

\begin{enumerate}
\def\labelenumi{(\alph{enumi})}
\setcounter{enumi}{1}
\item
  Let G be the free product of a family \((G_{\alpha})\) of groups. Show that, if all the \(G_{\alpha}\) are residually finite (cf.~§ 5, Exercise 5), so is G. (Reduce it first to the case where the family is finite, then to the case where it has two elements; by (a) it may be assumed that the \(G_{\alpha}\) are finite, then apply the preceding exercise.)
\item
  Deduce that every free group is residually finite.
\end{enumerate}

\begin{enumerate}
\def\labelenumi{\arabic{enumi}.}
\setcounter{enumi}{34}
\item
  *Let G (resp. G') be the additive group of fractions of the form \(a/b\) with \(a \in \mathbf{Z}\) , \(b \in \mathbf{Z}\) and \(b\) not divisible by 2 (resp. by 3). Let E = G * \(_{\mathbf{Z}}\) G' be the sum of G and G' amalgamated by their common subgroup Z. Show that G and G' are residually finite, but that E is not. (Prove that the only subgroup of E of finite index is E.)*
\item
  Let \((\mathbf{G}_1, \ldots, \mathbf{G}_n)\) and \((\mathbf{A}_1, \ldots, \mathbf{A}_{n-1})\) be two sequences of groups; for each \(i\) , let \(\phi_i: \mathbf{A}_i \to \mathbf{G}_i\) and \(\psi_i: \mathbf{A}_i \to \mathbf{G}_{i+1}\) be two injective homomorphisms; by these homomorphisms \(\mathbf{A}_i\) can be identified both with a subgroup of \(\mathbf{G}_i\) and with a subgroup of \(\mathbf{G}_{i+1}\) .
\end{enumerate}

\begin{enumerate}
\def\labelenumi{(\alph{enumi})}
\item
  Let \(H_{1}=G_{1}, H_{2}=H_{1} *_{A_{1}} G_{2}, \ldots, H_{n}=H_{n-1} *_{A_{n-1}} G_{n}\) . The group \(H_{n}\) is called the sum of the groups \((G_{i})\) amalgamated by the subgroups \((A_{i})\) ; it is denoted by \(G_{1} *_{A_{1}} G_{2} *_{A_{2}} \cdots *_{A_{n-1}} G_{n}\) . If T is an arbitrary group, define a bijection between the homomorphisms of \(H_{n}\) into T and the families \((t_{1},\ldots ,t_{n})\) of homomorphisms \(t_i\colon \mathbf{G}_i\to \mathbf{T}\) such that \(t_i\circ \phi_i = t_{i + 1}\circ \psi_i\) for \(1\leqslant i\leqslant n - 1\)
\item
  Each \(G_{i}\) is identified with the corresponding subgroup of \(H_{n}\) . Show that \(G_{i} \cap G_{i+1} = A_{i}\) and that \(G_{i} \cap G_{i+2} = A_{i} \cap A_{i+1}\) (the latter intersection being taken in \(G_{i+1}\) ).
\end{enumerate}

¶ 37. Let G be a group and S a finite symmetric subset of G generating G. Let \(\mathfrak{T}\) denote the set of finite subsets of G, ordered by inclusion; it is a directed set.

\begin{enumerate}
\def\labelenumi{(\alph{enumi})}
\item
  Let \(\mathbf{T} \in \mathfrak{T}\) and let \(\mathbf{B}_{\mathbf{T}}\) be the set of connected S-components of \(\mathbf{G} - \mathbf{T}\) (cf.~Exercise 18); let \(\mathbf{B}_{\mathbf{T}}^{\infty} = \mathbf{B}_{\mathbf{T}} - (\mathfrak{T} \cap \mathbf{B}_{\mathbf{T}})\) . Show that \(\mathbf{B}_{\mathbf{T}}\) is finite (note that, if \(\mathbf{T} \neq \varnothing\) and \(\mathbf{Z} \in \mathbf{B}_{\mathbf{T}}\) , there exists \(t \in \mathbf{T}\) and \(z \in \mathbf{Z}\) which are S-neighbours).
\item
  Let \(\mathbf{T}, \mathbf{T}' \in \mathfrak{T}\) with \(\mathbf{T} \subset \mathbf{T}'\) . If \(Z' \in \mathbf{B}_{\mathbf{T}'}\) , show that there exists a unique element \(Z \in \mathbf{B}_{\mathbf{T}}\) such that \(Z' \subset Z\) ; the mapping \(f_{\mathbf{T}\mathbf{T}'}: \mathbf{B}_{\mathbf{T}'} \to \mathbf{B}_{\mathbf{T}}\) such that \(f_{\mathbf{T}\mathbf{T}'}(\mathbf{Z}') = \mathbf{Z}\) maps \(\mathbf{B}_{\mathbf{T}}^{\infty}\) onto \(\mathbf{B}_{\mathbf{T}}^{\infty}\) . The inverse limit of the \(\mathbf{B}_{\mathbf{T}}\) , with \(\mathbf{T}\) running through \(\mathfrak{T}\) , is denoted by \(\mathbf{B}\) . Show that \(B = \varprojlim B_{\mathbf{T}}^{\infty}\) and that the image of \(B\) in \(B_{\mathbf{T}}\) is \(B_{\mathbf{T}}^{\infty}\) . The set \(B\) is called the set of ends of \(G\) ; it is non-empty if and only if \(G\) is infinite.
\item
  Let \(\mathbf{T}\) be a non-empty S-connected finite subset of \(\mathbf{G}\) and let \(\mathbf{Z} \in \mathbf{B}_{\mathbf{T}}^{\infty}\) . Show that there exist \(g \in \mathbf{G}\) such that \(g\mathbf{T} \cap \mathbf{T} = \varnothing\) and \(g\mathbf{T} \cap \mathbf{Z} \neq \varnothing\) and that then \(g\mathbf{T} \subset \mathbf{Z}\) . Show that there exists \(\mathbf{Z}_1 \in \mathbf{B}_{\mathbf{T}}\) such that \(g\mathbf{Z}_1\) contains \(\mathbf{T}\) and all the \(\mathbf{Z}' \in \mathbf{B}_{\mathbf{T}}\) such that \(\mathbf{Z}' \neq \mathbf{Z}\) (observe that \(\mathbf{T} \cup \bigcup_{\mathbf{Z}' \neq \mathbf{Z}} \mathbf{Z}'\) is S-connected and does not meet \(g\mathbf{T}\) ); deduce that, if \(\mathbf{Z}' \in \mathbf{B}_{\mathbf{T}}\) is different from \(\mathbf{Z}_1\) , then \(g\mathbf{Z}' \subset \mathbf{Z}\) . Show that, writing \(\mathbf{T}' = \mathbf{T} \cup g\mathbf{T}\) , the inverse image of \(\mathbf{Z}\) in \(\mathbf{B}_{\mathbf{T}}^{\infty}\) consists of at least \(n - 1\) elements, where \(n = \operatorname{Card}(\mathbf{B}_{\mathbf{T}}^{\infty})\) .
\item
  Let \(S'\) be a symmetric finite subset of \(G\) generating \(G\) and \(B'\) the set of ends of \(G\) relative to \(S'\) . Define a bijection of \(B\) onto \(B'\) (reduce it to the case where \(S \subset S'\) ). The cardinal of \(B\) is thus independent of the choice of \(S\) ; it is called the number of ends of \(G\) .
\item
  Show, using (c), that the number of ends of G is equal to 0, 1, 2 or \(2^{\aleph_0}\) .
\item
  Show that the number of ends of \(\mathbf{Z}\) is 2 and that that of \(\mathbf{Z}^n\) ( \(n \geqslant 2\) ) is 1.
\end{enumerate}

¶ 38. Let X be a finite set, F(X) the free group constructed on X and S = X ∪ X \(^{-1}\) .

\begin{enumerate}
\def\labelenumi{(\alph{enumi})}
\item
  Let \(\mathbf{S}_n\) be the set of products \(s_1 \ldots s_m\) , \(s_i \in \mathbf{S}\) , \(m \leqslant n\) . Show that every connected S-component (cf.~Exercise 18) of \(\mathrm{F}(\mathbf{X}) - \mathbf{S}_n\) contains one and only one element of the form \(s_1 \ldots s_{n+1}\) with \(s_i \in \mathbf{S}\) and \(s_i s_{i+1} \neq e\) for \(1 \leqslant i \leqslant n\) .
\item
  Deduce from (a) a bijection of the set B of ends of F(X) (cf.~Exercise 37) onto the set of infinite sequences \((s_{1},\ldots,s_{n},\ldots)\) of elements of S such that \(s_{i}s_{i+1}\neq e\) for all i. In particular, the number of ends of F(X) is equal to \(2^{\aleph_0}\) if \(\text{Card}(X)\geqslant2\) .
\end{enumerate}

\begin{enumerate}
\def\labelenumi{\arabic{enumi}.}
\setcounter{enumi}{38}
\tightlist
\item
  *Let G be a finitely generated group and F the set of symmetric finite subsets of G containing e, which generate G.
\end{enumerate}

\begin{enumerate}
\def\labelenumi{(\alph{enumi})}
\tightlist
\item
  If \(\mathbf{X} \in \mathfrak{F}\) and \(n \in \mathbf{N}\) , let \(d_n(\mathbf{X})\) denote the number of elements in \(\mathbf{G}\) of the form \(x_1 \ldots x_n\) with \(x_i \in \mathbf{X}\) . If \(\mathbf{X}, \mathbf{Y} \in \mathfrak{F}\) , show that there exists an integer \(a \geqslant 1\) such that
\end{enumerate}

\[
d _ {n} (\mathbf {X}) \leqslant d _ {a n} (\mathbf {Y}) \quad \text { for   all } n \in \mathbf {N}.
\]

Deduce that \(\limsup_{n\to \infty}\frac{\log(d_n(\mathbf{X}))}{\log(n)}\) does not depend on the choice of \(\mathbf{X}\) in \(\mathfrak{F}\) ; this limit (finite or infinite) is denoted by \(e(\mathbf{G})\) .

\begin{enumerate}
\def\labelenumi{(\alph{enumi})}
\setcounter{enumi}{1}
\item
  If \(\mathbf{X} \in \mathfrak{F}\) and \(\mathbf{G}\) is infinite, show that \(d_n(\mathbf{X}) < d_{n+1}(\mathbf{X})\) for all \(n\) . Deduce that \(e(\mathbf{G})\) is \(\geqslant 1\) .
\item
  Let H be a finitely generated group. Show that \(e(\mathbf{H}) \leqslant e(\mathbf{G})\) if H is isomorphic to a subgroup (resp. a quotient group) of G. If E is an extension of G by H, show that \(e(\mathbf{E}) \geqslant e(\mathbf{H}) + e(\mathbf{G})\) , with equality when \(E = G \times H\) .
\item
  Show that \(e(\mathbf{Z}^n) = n\) .
\item
  Let \(\mathbf{X} \in \mathfrak{F}\) . Consider the following property:
\item
  There exists a real number \(a > 1\) such that \(d_n(\mathbf{X}) \geqslant a^n\) for all sufficiently large \(n\) .
\end{enumerate}

Show that, if property (i) holds for one element of \(\mathfrak{F}\) , it holds for every element of \(\mathfrak{F}\) . \(\mathbf{G}\) is then said to be of exponential type; this implies \(e(\mathbf{G}) = +\infty\) .

\begin{enumerate}
\def\labelenumi{(\alph{enumi})}
\setcounter{enumi}{5}
\item
  Let \(G_{1}\) and \(G_{2}\) be two groups containing the same subgroup A and let \(H = G_{1} *_{A} G_{2}\) the corresponding amalgamated sum. Suppose \(G_{1}\) and \(G_{2}\) are finitely generated (hence also H is) and \(G_{i} \neq A\) for i = 1, 2. Show that H is of exponential type if at least one of the indices ( \(G_{1}:A\) ), ( \(G_{2}:A\) ) is \(\geqslant 3\) .
\item
  Show that a free group of rank \(\geqslant 2\) is of exponential type.*
\end{enumerate}

\begin{enumerate}
\def\labelenumi{\arabic{enumi}.}
\setcounter{enumi}{39}
\tightlist
\item
  *Let \(n\) be an integer \(\geqslant 2\) and \(T_n\) the group of upper triangular matrices of order \(n\) over \(Z\) with all the diagonal terms equal to 1. Let \(e(T_n)\) be the invariant of the group \(T_n\) defined in Exercise 39. Show that
\end{enumerate}

\[
e (\mathrm{T} _ {n}) \leqslant \sum_ {i = 1} ^ {n} i (n - i). *
\]

